% ============================================================
\chapter{Sectional, Ricci, and Scalar Curvature}
\label{ch:curvatures}
% ============================================================

The Riemann tensor is a rich but complex object---in dimension $n$, it has $\frac{n^2(n^2-1)}{12}$ independent components. To extract more readable geometric information, one contracts this tensor in different ways. \emph{Sectional curvature} measures curvature in each tangent plane---it is the direct generalisation of Gaussian curvature for surfaces. \emph{Ricci curvature}, obtained by averaging over planes containing a given direction, controls the convergence or divergence of geodesics and appears in Einstein's equations. \emph{Scalar curvature}, the final contraction, gives a single number at each point. These three levels of reading form the fundamental hierarchy of Riemannian geometry.

\section{Sectional curvature}

\begin{definition}[Sectional curvature]
Let $(M, g)$ be a Riemannian manifold and $\sigma \subset T_pM$ a $2$-plane
spanned by linearly independent vectors $u, v$. The \emph{sectional curvature}
of $\sigma$ is:
\[
    K(\sigma) = K(u, v) = \frac{\inner{R(u,v)v}{u}}{\norm{u}^2\norm{v}^2 - \inner{u}{v}^2}.
\]
\end{definition}

\begin{proposition}
The sectional curvature $K(u,v)$ depends only on the $2$-plane
$\sigma = \operatorname{Span}(u,v)$, not on the choice of basis.
\end{proposition}

\begin{proof}
If $(u', v') = (au + bv, cu + dv)$ is another basis of $\sigma$ with $ad - bc \neq 0$,
then by the symmetries of $R$:
$\inner{R(u',v')v'}{u'} = (ad-bc)^2 \inner{R(u,v)v}{u}$
and $\norm{u'}^2\norm{v'}^2 - \inner{u'}{v'}^2 = (ad-bc)^2(\norm{u}^2\norm{v}^2 - \inner{u}{v}^2)$.
The ratio is therefore invariant.
\end{proof}

\begin{theorem}[Determination of the Riemann tensor]
The sectional curvature completely determines the Riemann tensor. More precisely,
if $K(u,v)$ is known for all $2$-planes, then $R$ is uniquely determined.
\end{theorem}

\begin{proof}
The function $(u,v) \mapsto \inner{R(u,v)v}{u}$ is a quadratic form in each
variable $u$ and $v$. By the symmetries of $R$, knowing this form for all
$u, v$ determines $R_{ijkl}$ by polarization.
\end{proof}

\section{Spaces of constant curvature}

\begin{definition}
$(M, g)$ is a \emph{space of constant (sectional) curvature} $\kappa$ if
$K(\sigma) = \kappa$ for every $2$-plane $\sigma$ at every point.
\end{definition}

\begin{theorem}[Form of the Riemann tensor]
If $K \equiv \kappa$, then:
\[
    R(X,Y)Z = \kappa\big(\inner{Y}{Z}X - \inner{X}{Z}Y\big).
\]
In coordinates: $R_{ijkl} = \kappa(g_{ik}g_{jl} - g_{il}g_{jk})$.
\end{theorem}

\begin{keyformulas}[title=\textbf{Spaces of constant curvature}]
\begin{center}
\renewcommand{\arraystretch}{1.3}
\begin{tabular}{lccc}
\hline
\textbf{Space} & $\kappa$ & \textbf{Model} & \textbf{Topology} \\
\hline
Sphere $S^n$ & $+1$ & $\{x \in \R^{n+1} : \norm{x} = 1\}$ & compact, $\pi_1 = 0$ \\
Euclidean $\R^n$ & $0$ & $\R^n$ & noncompact, $\pi_1 = 0$ \\
Hyperbolic $\mathbb{H}^n$ & $-1$ & half-space & noncompact, $\pi_1 = 0$ \\
\hline
\end{tabular}
\end{center}
\end{keyformulas}

\begin{theorem}[Classification -- Killing--Hopf]
Let $(M, g)$ be a complete simply connected Riemannian manifold of constant
sectional curvature $\kappa$. Then $(M, g)$ is isometric to:
\begin{itemize}
    \item $S^n(1/\sqrt{\kappa})$ if $\kappa > 0$,
    \item $\R^n$ if $\kappa = 0$,
    \item $\mathbb{H}^n(1/\sqrt{-\kappa})$ if $\kappa < 0$.
\end{itemize}
Every complete manifold of constant curvature is a quotient of these spaces by a
discrete group of isometries acting freely and properly.
\end{theorem}

\section{Ricci curvature}

\begin{definition}[Ricci tensor]
The \emph{Ricci tensor} is the contraction of the Riemann tensor:
\[
    \operatorname{Ric}(X, Y) = \tr\big(Z \mapsto R(Z, X)Y\big).
\]
In coordinates: $R_{ij} = R^k_{ikj} = g^{kl} R_{kilj}$.
\end{definition}

\begin{proposition}
The Ricci tensor is symmetric: $\operatorname{Ric}(X,Y) = \operatorname{Ric}(Y,X)$.
\end{proposition}

\begin{proof}
In an orthonormal basis:
$R_{ij} = \sum_k R_{kikj}$ and $R_{ji} = \sum_k R_{kjki} = \sum_k R_{kikj} = R_{ij}$,
using the pair symmetry $R_{kjki} = R_{kikj}$.
\end{proof}

\begin{proposition}[Geometric interpretation]
If $(e_1, \ldots, e_n)$ is an orthonormal basis of $T_pM$ with $e_1 = v/\norm{v}$:
\[
    \operatorname{Ric}(v, v) = \norm{v}^2 \sum_{i=2}^n K(v, e_i).
\]
The Ricci curvature in direction $v$ is the \emph{sum of sectional curvatures}
of $2$-planes containing $v$.
\end{proposition}

\begin{example}[Ricci of $S^n$]
All sectional curvatures equal $1$, so
$\operatorname{Ric}(v,v) = (n-1)\norm{v}^2$, i.e.\ $\operatorname{Ric} = (n-1)g$.
\end{example}

\begin{example}[Ricci of $\mathbb{H}^n$]
$\operatorname{Ric} = -(n-1)g$.
\end{example}

\begin{example}[Ricci of a bi-invariant Lie group]
With $R(X,Y)Z = \frac{1}{4}[[X,Y],Z]$:
\[
    \operatorname{Ric}(X,Y) = -\frac{1}{4}\tr(\operatorname{ad}_X \circ \operatorname{ad}_Y)
    = -\frac{1}{4}B(X,Y),
\]
where $B$ is the Killing form of $\mathfrak{g}$.
\end{example}

\section{Scalar curvature}

\begin{definition}[Scalar curvature]
The \emph{scalar curvature} is the trace of the Ricci tensor:
\[
    S = \operatorname{Scal} = \tr_g \operatorname{Ric} = g^{ij} R_{ij}.
\]
If $(e_1, \ldots, e_n)$ is orthonormal:
$S = \sum_{i} \operatorname{Ric}(e_i, e_i) = 2\sum_{i < j} K(e_i, e_j)$.
\end{definition}

\begin{example}
For $S^n$: $S = n(n-1)$. For $\mathbb{H}^n$: $S = -n(n-1)$. For $\R^n$: $S = 0$.
\end{example}

\section{Einstein manifolds}

\begin{definition}[Einstein manifold]
$(M, g)$ is an \emph{Einstein manifold} if the Ricci tensor is proportional to
the metric:
\[
    \operatorname{Ric} = \lambda\, g
\]
for a constant $\lambda \in \R$. Taking the trace: $S = n\lambda$, so $\lambda = S/n$.
\end{definition}

\begin{example}
Spaces of constant curvature are Einstein manifolds with $\lambda = (n-1)\kappa$.
\end{example}

\begin{theorem}[Schur]
If $n \geq 3$ and the sectional curvature is constant at each point (i.e.\
$K_p(\sigma) = f(p)$ for all $2$-planes $\sigma$ at $p$), then $f$ is constant
on $M$.

More generally, if $n \geq 3$ and $\operatorname{Ric} = f\, g$ with
$f \in C^\infty(M)$, then $f$ is constant.
\end{theorem}

\begin{proof}
From $\operatorname{Ric} = f\, g$ we get $S = nf$, and the contracted second
Bianchi identity gives:
\[
    \nabla^j R_{ij} = \frac{1}{2}\nabla_i S.
\]
But $\nabla^j R_{ij} = \nabla^j(fg_{ij}) = \nabla_i f$. Hence
$\nabla_i f = \frac{n}{2}\nabla_i f$, giving $(1 - n/2)\nabla_i f = 0$.
For $n \geq 3$, this implies $\nabla f = 0$.
\end{proof}

\section{The Weyl tensor}

\begin{definition}[Weyl tensor]
For $n \geq 3$, the \emph{Weyl tensor} $W$ is the part of the Riemann tensor
not determined by Ricci:
\begin{align*}
    R_{ijkl} &= W_{ijkl} + \frac{1}{n-2}\left(R_{ik}g_{jl} - R_{il}g_{jk}
    + R_{jl}g_{ik} - R_{jk}g_{il}\right) \\
    &\quad - \frac{S}{(n-1)(n-2)}\left(g_{ik}g_{jl} - g_{il}g_{jk}\right).
\end{align*}
\end{definition}

\begin{proposition}
The Weyl tensor is conformally invariant: if $\tilde{g} = e^{2\varphi}g$, then
$\tilde{W}^i_{jkl} = W^i_{jkl}$.
\end{proposition}

\begin{proposition}
$W \equiv 0$ if and only if $n \leq 3$, or if $n \geq 4$ and $(M,g)$ is
\emph{conformally flat} (locally conformal to $\R^n$).
\end{proposition}

\section{Curvature of warped products}

\begin{proposition}[Warped product curvature]
On $B \times_f F$ with $g = g_B + f^2 g_F$, when $\dim B = 1$ (i.e.\
$B = I \subset \R$, $g_B = dt^2$):
\begin{enumerate}
    \item For $V, W$ tangent to $F$:
    $K(V, W) = \frac{K_F(V,W) - (f')^2}{f^2}$.
    \item For $V$ tangent to $F$ and $\partial_t$ tangent to $B$:
    $K(\partial_t, V) = -\frac{f''}{f}$.
\end{enumerate}
\end{proposition}

\begin{example}[Curvature of $\R^n$ in polar coordinates]
$g = dr^2 + r^2 g_{S^{n-1}}$, so $f(r) = r$, $f' = 1$, $f'' = 0$.
$K(\partial_r, V) = 0$ and $K(V, W) = \frac{1 - 1}{r^2} = 0$. Everything is flat.
\end{example}

\begin{example}[Curvature of $S^n$ in geodesic coordinates]
$g = dr^2 + \sin^2\!r\, g_{S^{n-1}}$, $f(r) = \sin r$, $f'' = -\sin r$.
$K(\partial_r, V) = 1$ and $K(V,W) = \frac{1 - \cos^2 r}{\sin^2 r} = 1$.
\end{example}

\begin{example}[Curvature of $\mathbb{H}^n$ in geodesic coordinates]
$g = dr^2 + \sinh^2\!r\, g_{S^{n-1}}$, $f(r) = \sinh r$, $f'' = \sinh r$.
$K(\partial_r, V) = -1$ and $K(V,W) = \frac{1 - \cosh^2 r}{\sinh^2 r} = -1$.
\end{example}

\section{Contracted Bianchi identity}

\begin{theorem}[Contracted Bianchi identity]
\[
    \nabla^j R_{ij} = \frac{1}{2}\nabla_i S,
\]
or equivalently, the \emph{Einstein tensor}
$G_{ij} = R_{ij} - \frac{1}{2}Sg_{ij}$ is divergence-free:
$\nabla^j G_{ij} = 0$.
\end{theorem}

\begin{remark}
This identity is fundamental in general relativity, where Einstein's equations
read $G_{ij} = 8\pi T_{ij}$, and energy-momentum conservation
$\nabla^j T_{ij} = 0$ is automatic.
\end{remark}

\section{Exercises}

\begin{exercise}
Compute the sectional curvature, Ricci tensor, and scalar curvature of the torus
$S^1(r_1) \times S^1(r_2)$ with the product metric.
\end{exercise}

\begin{exercise}
Show that if $(M, g)$ has dimension $2$, then $\operatorname{Ric} = K\, g$ and $S = 2K$.
\end{exercise}

\begin{exercise}
Show that $\C P^n$ with the Fubini--Study metric is an Einstein manifold with
$\lambda = 2(n+1)$.
\end{exercise}

\begin{exercise}
Compute the Weyl tensor of $S^2 \times S^2$ and show it is not conformally flat
for $n = 4$.
\end{exercise}

\begin{exercise}
For the warped product $g = dt^2 + e^{2t} g_{\R^{n-1}}$, compute all sectional
curvatures, the Ricci tensor, and the scalar curvature.
\end{exercise}

\begin{exercise}
Prove Schur's theorem in dimension $3$ using the second Bianchi identity. Give
a counterexample showing the theorem fails in dimension $2$.
\end{exercise}
